% from ega3-1-fr.tex
\oldpage[0\textsubscript{III}]{32}
    ou pour $j<j_0$).
  \item[(b)] Il existe $i_0$ et $i_1$ tels que $K_{ij}=0$ pour $i<i_0$ et
    $i>i_1$.
  \item[(c)] Il existe $j_0$ et $j_1$ tels que $K_{ij}=0$ pour $j<j_0$ et
    $j>j_1$.
\end{enumerate}

La suite $'E(K_{\bullet\bullet})$ est régulière s'il existe $i_0$ tel que
$K_{ij}=0$ pour $i<i_0$, ou s'il existe $j_0$ tel que $K_{ij}=0$ pour
$j>j_0$.

La suite $''E(K_{\bullet\bullet})$ est régulière s'il existe $i_0$ tel que
$K_{ij}=0$ pour $i>i_0$, ou s'il existe $j_0$ tel que $K_{ij}=0$ pour
$j<j_0$.
\end{env}

\subsection{Hypercohomologie d'un foncteur par rapport à un complexe
$K^\bullet$}
\label{subsection:0.11.4-fr}

\begin{env}[11.4.1]
\label{0.11.4.1-fr}
Soit $\mathcal C$ une catégorie abélienne~; rappelons que l'on appelle
\emph{résolution droite} (ou \emph{cohomologique}) d'un objet $A$ de
$\mathcal C$ un complexe d'objets de $\mathcal C$, dont l'opérateur de
dérivation est de degré $+1$,
\[
  0\longrightarrow L^0\longrightarrow L^1\longrightarrow L^2
  \longrightarrow\cdots
\]
muni d'un morphisme $\varepsilon:A\to L^0$ dit \emph{augmentation} de la
résolution (et que l'on peut considérer comme un morphisme de complexes
\[
\xymatrix{
  0\ar[r] & A\ar[r]\ar[d]_{\varepsilon} & 0\ar[r]\ar[d]
    & 0\ar[r]\ar[d] & \cdots\\
  0\ar[r] & L^0\ar[r] & L^1\ar[r] & L^2\ar[r] & \cdots
})
\]
tel que la suite
\[
  0\longrightarrow A\xrightarrow{\varepsilon}L^0\longrightarrow L^1
  \longrightarrow\cdots
\]
soit exacte~; de même une \emph{résolution gauche} (ou \emph{homologique})
de $A$ est un complexe
$0\leftarrow L_0\leftarrow L_1\leftarrow\cdots$ d'objets de $\mathcal C$
dont l'opérateur de dérivation est de degré $-1$, muni d'une augmentation
$\varepsilon:L_0\to A$, de sorte que la suite
\[
  0\longleftarrow A\xleftarrow{\varepsilon}L_0\longleftarrow L_1
  \longleftarrow\cdots
\]
soit exacte.

Lorsqu'une résolution droite $(L_i)_{i\geq0}$ d'un objet $A$ est telle que
$L_i=0$ pour $i\geq n+1$, on dit que cette résolution est de
\emph{longueur $\leq n$}. On définit de même une résolution gauche de
longueur $\leq n$. Une résolution qui est de longueur $\leq n$ pour un entier
$n$ est dite \emph{finie}.

Une résolution de $A$ est dite \emph{projective} (resp. \emph{injective}) si
les objets de $\mathcal C$ autres que $A$ qui la composent sont
\emph{projectifs} (resp. \emph{injectifs}). Lorsque $\mathcal C$ est la
catégorie des modules (à gauche par exemple) sur un anneau, on dira de même
qu'une résolution de $A$ est \emph{plate} (resp. \emph{libre}) lorsque les
modules autres que $A$ qui la composent sont \emph{plats} (resp.
\emph{libres}).
\end{env}

\begin{env}[11.4.2]
\label{0.11.4.2-fr}
Soit $K^\bullet=(K^i)_{i\in\mathbb Z}$ un complexe d'objets de $\mathcal C$,
dont l'opérateur de dérivation est de degré $+1$.

On appelle \emph{résolution de Cartan-Eilenberg droite} de $K^\bullet$ le
couple formé d'un bicomplexe $L^{\bullet\bullet}=(L^{ij})$ à opérateurs de
dérivation de degré $+1$, avec $L^{ij}=0$ pour $j<0$, et d'un morphisme de
complexes simples $\varepsilon:K^\bullet\to L^{\bullet,0}$, de façon que les
conditions suivantes soient remplies~:
